\documentclass[11pt]{article}
\usepackage[margin=1in]{geometry}
\usepackage{amsmath,amssymb,amsthm}
\usepackage{xcolor}
\usepackage[colorlinks=true,linkcolor=blue!55!black,citecolor=blue!55!black,
urlcolor=blue!65!black]{hyperref}
\usepackage{microtype}

\newtheorem{proposition}{Proposition}
\newcommand{\R}{\mathbb R}

\title{\textcolor{blue!55!black}{Fisher's inequality excludes two proposed
Steiner ETF candidates}}
\author{Automated mathematics run \texttt{fa\_banach\_001}, lane 13}
\date{23 July 2026}

\begin{document}
\maketitle

\noindent\fbox{\parbox{0.94\linewidth}{
\textbf{Status.} Literature-implied answer, complete for the two finite
candidates.  The broader ETF restricted-isometry question remains open.
}}

\section{Source question}

Fickus, Mixon, and Tremain \cite{FMT}, at the end of Section 3 on page 14,
write that they were unable to resolve two remaining candidates:
the \(23\times276\) and \(46\times736\) real equiangular tight frames could
potentially arise from \((2,10,46)\)- and \((2,14,92)\)-Steiner systems,
respectively, provided those systems exist.

\section{Identification}

The classical Fisher inequality for block designs settles both cases.
Horsley \cite[p.~2]{Horsley} recalls that every nontrivial
\(2\)-\((v,k,\lambda)\) design has at least \(v\) blocks, equivalently
\[
 \lambda(v-1)\geq k(k-1).
\]

\begin{proposition}
Neither a \((2,10,46)\)-Steiner system nor a
\((2,14,92)\)-Steiner system exists.  Consequently, neither of the two
source ETFs can arise from the proposed Steiner construction.
\end{proposition}

\begin{proof}
Let \(A\) be the \(v\times b\) incidence matrix of a nontrivial
\(2\)-\((v,k,\lambda)\) design, and let \(r\) be its replication number.
Then
\[
 AA^{\mathsf T}=(r-\lambda)I+\lambda J.
\]
Because \(v>k\), one has \(r>\lambda\).  The displayed matrix is therefore
positive definite: its eigenvalues are \(r-\lambda\), with multiplicity
\(v-1\), and \(r-\lambda+\lambda v\), with multiplicity one.  Hence
\[
 v=\operatorname{rank}(AA^{\mathsf T})
 \leq \operatorname{rank}(A)\leq b,
\]
which is Fisher's inequality.

For a \((2,k,v)\)-Steiner system, \(\lambda=1\) and double counting gives
\[
 b=\frac{v(v-1)}{k(k-1)}.
\]
The first proposed system would have
\[
 b=\frac{46\cdot45}{10\cdot9}=23<46,
\]
and the second would have
\[
 b=\frac{92\cdot91}{14\cdot13}=46<92.
\]
Both contradict \(b\geq v\).
\end{proof}

\section{Scope and provenance}

The argument gives a complete negative answer to the two explicit finite
candidates and also records the general necessary condition
\[
 v-1\geq k(k-1)
\]
for every nontrivial \((2,k,v)\)-Steiner system.  This is not a new theorem:
it is a direct application of Fisher's classical inequality, apparently
missed in the source discussion.  It is therefore classified as a
literature-implied answer rather than as a new full solution.

The source's final question---whether some other equiangular tight frames
can have restricted-isometry behavior beyond the square-root bottleneck---is
not answered here.

\begin{thebibliography}{9}
\bibitem{FMT}
M.~Fickus, D.~G.~Mixon, and J.~C.~Tremain,
\emph{Steiner equiangular tight frames},
Linear Algebra Appl. \textbf{436} (2012), 1014--1027;
\href{https://arxiv.org/abs/1009.5730}{arXiv:1009.5730}.

\bibitem{Horsley}
D.~Horsley,
\emph{Generalising Fisher's inequality to coverings and packings},
Australas. J. Combin. \textbf{70} (2018), 339--372;
\href{https://arxiv.org/abs/1409.0485}{arXiv:1409.0485}.
\end{thebibliography}

\end{document}
